\section{Background and Related Work}
\label{sec:background}

\subsection{Vehicle Routing with Synchronization Constraints}
Home healthcare routing and scheduling has emerged as a major optimization challenge due to population aging, the increasing prevalence of chronic diseases, and the growing adoption of decentralized healthcare services \citep{masmoudi2024home}. In such environments, caregivers must travel between geographically dispersed patients while respecting operational constraints including service time windows, staff qualifications, workload balancing, and route duration limitations.

Many home healthcare activities additionally require coordinated interventions involving multiple caregivers. Examples include treatments requiring complementary medical skills, patient transfers, rehabilitation sessions, and procedures involving specialized equipment. These situations introduce synchronization constraints, creating temporal dependencies between visits that must be performed simultaneously or according to predefined precedence relationships.

The Vehicle Routing Problem with Time Windows and Synchronization constraints (VRPTW-S) models these operational requirements by combining routing, scheduling, and synchronization decisions within a unified optimization framework. Compared with the classical Vehicle Routing Problem with Time Windows (VRPTW), synchronization constraints significantly increase problem complexity by coupling routes that would otherwise be independent. Consequently, modifications performed on a given route may propagate temporal infeasibilities to several other routes simultaneously, making feasibility management considerably more difficult.

The VRPTW-S belongs to the class of NP-hard combinatorial optimization problems \citep{lenstra1981complexity}. Exact optimization methods rapidly become computationally expensive as instance size and synchronization density increase, motivating the development of advanced heuristic and metaheuristic approaches capable of producing high-quality solutions within practical computational times.

\subsection{Metaheuristic Approaches for VRPTW-S}
Early research on the VRPTW-S primarily focused on exact optimization techniques. Among the pioneering contributions, \citet{bredstrom2008combined} proposed a branch-and-price framework for synchronized home healthcare routing and introduced benchmark instances that subsequently became standard references in the literature. While exact methods can provide optimal solutions for small and medium-sized instances, their scalability remains limited due to the strong temporal dependencies induced by synchronization constraints.

As problem sizes increased and synchronization requirements became more prevalent in realistic healthcare applications, research progressively shifted toward heuristic and metaheuristic approaches capable of balancing solution quality and computational effort. Among the earliest successful methods, \citet{afifi2016heuristic} developed a synchronization-aware Simulated Annealing framework incorporating specialized neighborhood operators and adaptive penalty mechanisms. Their results highlighted the importance of explicitly considering synchronization constraints during local search exploration.

Subsequent work further improved solution quality through adaptive search mechanisms. \citet{liu2019adaptive} proposed an Adaptive Large Neighborhood Search (ALNS) algorithm combining destroy-and-repair operators specifically designed for synchronized routing. By dynamically adjusting operator scores according to their historical performance, the method achieved several improvements over previously reported benchmark results.

Among neighborhood-based metaheuristics, Variable Neighborhood Search (VNS) and its extensions have demonstrated particularly strong performance. The Generalized Variable Neighborhood Search (GVNS) introduced by \citet{masmoudi2025generalized} currently represents one of the most competitive approaches reported for the VRPTW-S. Their framework combines shaking procedures, Variable Neighborhood Descent (VND), and adaptive penalty adjustment mechanisms to efficiently explore both feasible and infeasible regions of the search space. Extensive computational experiments reported numerous best-known solutions across standard benchmark sets, establishing GVNS as a strong reference method for synchronization-constrained routing optimization.

Despite the significant progress achieved by modern metaheuristics, most existing approaches continue to rely on search-control decisions that remain fixed throughout optimization. Neighborhood exploration orders, local search strategies, acceptance mechanisms, and algorithmic parameters are generally determined offline and remain unchanged regardless of the current search state. Such static decisions may limit adaptability to the highly dynamic search landscape induced by synchronization constraints.

\subsection{Learning-Guided Optimization for Routing Problems}
Recent advances in machine learning have stimulated growing interest in learning-guided combinatorial optimization. Rather than relying exclusively on handcrafted search policies, these approaches seek to learn adaptive decision mechanisms capable of exploiting information gathered during the optimization process. Comprehensive surveys of this emerging field are provided by \citet{karimi2022machine}, \citet{mazyavkina2021reinforcement}, and \citet{bengio2021machine}.

Early studies primarily focused on learning complete solution construction policies. \citet{bello2017neural} demonstrated that reinforcement learning could be used to generate solutions for the Traveling Salesman Problem without explicit supervision. This idea was subsequently extended to vehicle routing problems by \citet{nazari2018reinforcement}, who proposed a reinforcement learning framework for the Capacitated Vehicle Routing Problem, and by \citet{kool2019attention}, who introduced attention-based architectures capable of producing competitive routing solutions across several combinatorial optimization tasks.

More recently, learning has increasingly been used not to replace optimization algorithms, but to guide their search behavior. Existing approaches can broadly be categorized according to the level at which learning influences the optimization process. Some methods focus on adaptive operator selection, others target parameter control, move evaluation, acceptance decisions, or search strategy adaptation. This direction is particularly attractive for complex constrained routing problems because it preserves the strengths of established metaheuristics while introducing adaptive decision-making capabilities.

Several studies have demonstrated the effectiveness of learning-guided search control in routing optimization. For example, \citet{xu2023learning} proposed a reinforcement learning framework for adaptive neighborhood selection within a Variable Neighborhood Search procedure, showing that search performance can be improved through context-aware operator selection. Reinforcement learning has also been investigated for dynamic parameter adaptation, where policy-gradient methods and Proximal Policy Optimization (PPO) have shown promising results for controlling search behavior in complex optimization environments \citep{schulman2017proximal, wang2024csppo}. Similarly, supervised learning models have been employed to estimate move quality, predict operator effectiveness, and support adaptive search decisions during optimization.

Although these studies demonstrate the potential of learning-guided optimization, the vast majority have focused on classical routing variants that do not include synchronization constraints. Research investigating learning-based search adaptation for synchronization-constrained vehicle routing remains extremely limited, and the interaction between learning mechanisms and synchronization-aware metaheuristics remains largely unexplored.

\subsection{Research Gap and Contributions}
The literature highlights two complementary observations. First, state-of-the-art approaches for the VRPTW-S are predominantly based on advanced metaheuristics, with GVNS currently representing one of the most effective frameworks reported in the literature \citep{masmoudi2025generalized}. Second, recent studies in combinatorial optimization have demonstrated that learning-guided search control can successfully improve decision making across several routing variants.

However, these two research directions have remained largely disconnected. Existing VRPTW-S approaches continue to rely on static search-control decisions at multiple levels of the optimization process. Neighborhood operators are explored according to predefined rules, local search strategies remain fixed throughout execution, and algorithm parameters are typically tuned offline. Whether adaptive learning mechanisms can effectively replace these static decisions in synchronization-constrained routing remains largely unanswered.

To the best of our knowledge, learning-based search adaptation has received very limited attention in the VRPTW-S literature, and no previous work has investigated the combined integration of reinforcement learning and supervised learning mechanisms within a synchronization-aware GVNS framework.

This work addresses this gap through a two-stage contribution. First, we introduce GVNS+, an enhanced implementation of the state-of-the-art GVNS framework incorporating additional neighborhood operators and a new family of large-scale synchronized benchmark instances extending existing literature datasets up to 400 customers. Since existing synchronized benchmark sets remain relatively limited in scale, these larger instances were generated to enable a more comprehensive assessment of scalability, robustness, and learning generalization. The enhanced implementation is validated against the benchmark results reported by \citet{masmoudi2025generalized}, achieving an average gap of 1.14\% despite being entirely reimplemented in Python.

Second, we investigate the impact of learning-guided search adaptation on top of this enhanced baseline. Three complementary adaptive mechanisms are studied: reinforcement learning for neighborhood operator selection, supervised learning for dynamic First Improvement versus Best Improvement strategy selection, and Proximal Policy Optimization for online parameter control. Beyond evaluating the effectiveness of individual learning techniques, the study examines how the benefits of learning vary across different decision layers of the metaheuristic search process, providing insights into where adaptive intelligence can most effectively complement synchronization-aware optimization.